\section{연습문제}
\label{sec:separability-exercises}

문제는 A, B, C의 세 단계로 나뉜다. 별표 문제는 다음 두 절에 상세 풀이가 있다.

\subsection*{A. 판정과 계산}

\begin{exercise}[*]
다음 다항식이 주어진 체 위에서 분리다항식인지 판정하고, 비분리인 경우 중근의 중복도를 구하라.
\begin{enumerate}[label=(\alph*)]
  \item $X^4-2\in\Q[X]$,
  \item $X^4+X^2+1\in\F_2[X]$,
  \item $X^6-X\in\F_3[X]$.
\end{enumerate}
\end{exercise}

\begin{exercise}[*]
$K=\F_p(t)$에서 $f=X^p-t$를 생각하자.
\begin{enumerate}[label=(\alph*)]
  \item $f$가 $K[X]$에서 기약임을 보여라.
  \item $f$의 분리차수와 한 근을 첨가한 확장의 차수를 구하라.
  \item 대수적 폐포에서 $f$의 인수분해를 적어라.
\end{enumerate}
\end{exercise}

\begin{exercise}[*]
$L=\Q(\sqrt2,\sqrt3)$의 모든 $\Q$-매장 $L\to\C$를 구하고
\[
  [L:\Q]_{\mathrm s}
\]
를 계산하라. 각 매장이 $\sqrt2+\sqrt3$을 어디로 보내는지도 적어라.
\end{exercise}

\begin{exercise}
$\alpha=\sqrt[3]{2}$라 하자. $\Q(\alpha)$에서 $\C$로 가는 $\Q$-매장을 모두 구하라. 이 확장이 분리이지만 분해체는 아님을 설명하라.
\end{exercise}

\begin{exercise}
$\alpha^3+\alpha+1=0$인 $\alpha$에 대해 $L=\F_2(\alpha)$라 하자. 최소다항식이 분리임을 확인하고 $[L:\F_2]_{\mathrm s}$를 구하라.
\end{exercise}

\begin{exercise}[*]
$\theta=\sqrt2+\sqrt3$라 하자.
\begin{enumerate}[label=(\alph*)]
  \item $\sqrt2,\sqrt3\in\Q(\theta)$임을 보여라.
  \item $\theta$의 $\Q$ 위 최소다항식을 구하라.
  \item $\Q(\sqrt2,\sqrt3)=\Q(\theta)$임을 차수로도 확인하라.
\end{enumerate}
\end{exercise}

\begin{exercise}
$K=\F_p(t^p)$, $L=\F_p(t)$라 하자. $[L:K]$, $[L:K]_{\mathrm s}$와 비분리차수를 각각 구하라.
\end{exercise}

\begin{exercise}[*]
독립인 변수 $s,t$에 대해
\[
  K=\F_p(s^p,t^p),
  \qquad
  L=\F_p(s,t)
\]
라 하자.
\begin{enumerate}[label=(\alph*)]
  \item $1,s,\dots,s^{p-1}$가 $K(t)$ 위의 기저가 됨을 보여라.
  \item $\{s^it^j:0\le i,j<p\}$가 $L/K$의 기저임을 보여라.
  \item 보통 차수, 분리차수와 비분리차수를 구하라.
\end{enumerate}
\end{exercise}

\begin{exercise}
다음 체가 완전체인지 판정하라.
\[
  \Q,\quad \R,\quad \C,\quad \F_p,
  \quad \F_p(t),\quad \overline{\F_p}.
\]
\end{exercise}

\begin{exercise}[*]
$p$가 홀수이고 $s,t$가 독립인 변수라 하자. $K=\F_p(s,t)$에서
\[
  \alpha^p=s,
  \qquad
  \beta^2=t,
  \qquad
  L=K(\alpha,\beta)
\]
라 하자. $[L:K]$, $[L:K]_{\mathrm s}$와 $[L:K]_{\mathrm i}$를 구하고, $L$ 안에서의 $K$의 분리폐포를 추측하라.
\end{exercise}

\subsection*{B. 핵심 정리의 증명}

\begin{exercise}[*]
표수 $p>0$인 체 $K$와 $f\in K[X]$에 대해
\[
  f'=0
  \quad\Longleftrightarrow\quad
  f(X)=g(X^p)\text{인 }g\in K[X]\text{가 존재}
\]
함을 증명하라.
\end{exercise}

\begin{exercise}
표수 $p>0$에서 기약다항식 $f$의 차수가 $p$로 나누어지지 않으면 $f$가 분리다항식임을 증명하라.
\end{exercise}

\begin{exercise}[*]
$\charac K=p>0$이고 $\alpha$가 $K$ 위에서 대수적이라 하자. 다음 동치를 증명하라.
\[
  \alpha\text{가 }K\text{ 위에서 분리적}
  \quad\Longleftrightarrow\quad
  K(\alpha)=K(\alpha^p).
\]
\end{exercise}

\begin{exercise}[*]
유한 대수적 확장탑 $K\subseteq L\subseteq M$에 대해
\[
  [M:K]_{\mathrm s}
  =[M:L]_{\mathrm s}[L:K]_{\mathrm s}
\]
를 매장의 제한 사상과 연장 정리를 이용하여 증명하라.
\end{exercise}

\begin{exercise}
유한 대수적 확장 $L=K(\alpha_1,\dots,\alpha_n)$에서 모든 $\alpha_i$가 $K$ 위에서 분리적이면 $L/K$가 분리확장임을 증명하라.
\end{exercise}

\begin{exercise}
대수적 확장탑 $K\subseteq L\subseteq M$에서 $M/L$과 $L/K$가 분리확장이면 $M/K$도 분리확장임을 증명하라.
\end{exercise}

\begin{exercise}
체의 곱셈군의 모든 유한부분군이 순환군임을 증명하라.
\end{exercise}

\begin{exercise}[*]
$K$가 무한체이고 $L=K(\alpha,\beta)$가 유한 분리확장이라 하자. 유한개의 금지된 $c\in K$를 제외하면
\[
  L=K(\alpha+c\beta)
\]
임을 증명하라.
\end{exercise}

\begin{exercise}[*]
표수 $p>0$인 체 $K$에 대해 다음 동치를 증명하라.
\[
  K\text{가 완전체}
  \quad\Longleftrightarrow\quad
  K\to K,\ a\mapsto a^p\text{가 전사}.
\]
\end{exercise}

\begin{exercise}
대수적 확장 $L/K$에서 $K$ 위에서 분리적인 모든 원소의 집합이 중간체를 이룸을 증명하라.
\end{exercise}

\subsection*{C. 구조와 응용}

\begin{exercise}[*]
$K=\F_p(s^p,t^p)$, $L=\F_p(s,t)$라 하자. 모든 $\gamma\in L$에 대해 $\gamma^p\in K$임을 이용하여 $L/K$가 단순확장이 아님을 증명하라.
\end{exercise}

\begin{exercise}
대수적 확장 $L/K$에 대해 다음 조건들의 동치성을 증명하라.
\begin{enumerate}[label=(\roman*)]
  \item $L/K$는 순수비분리 확장이다.
  \item 모든 $\alpha\in L$에 대해 어떤 $r\ge0$이 존재하여 $\alpha^{p^r}\in K$이다.
  \item $K$를 고정하는 $L\to\overline K$ 매장이 하나뿐이다.
\end{enumerate}
\end{exercise}

\begin{exercise}
순수비분리성의 추이성을 $p$-거듭제곱 조건만을 이용하여 증명하라.
\end{exercise}

\begin{exercise}[*]
$L/K$가 대수적 확장이고 $\alpha,\beta\in L$라 하자. $\alpha$는 $K$ 위에서 분리적이고 $\beta$는 $K$ 위에서 순수비분리적이라고 하자.
\begin{enumerate}[label=(\alph*)]
  \item $K(\alpha,\beta)=K(\alpha+\beta)$임을 증명하라.
  \item $\alpha\beta\ne0$이면 $K(\alpha,\beta)=K(\alpha\beta)$임을 증명하라.
\end{enumerate}
\end{exercise}

\begin{exercise}[*]
대수적 확장 $L/K$에서
\[
  K_{\mathrm s}=\{\alpha\in L:\alpha\text{가 }K\text{ 위에서 분리적}\}
\]
라 하자. $K_{\mathrm s}/K$가 분리이고 $L/K_{\mathrm s}$가 순수비분리이며, 이 성질을 갖는 중간체가 유일함을 증명하라.
\end{exercise}

\begin{exercise}
유한 대수적 확장탑 $K\subseteq L\subseteq M$에 대해 비분리차수가
\[
  [M:K]_{\mathrm i}
  =[M:L]_{\mathrm i}[L:K]_{\mathrm i}
\]
를 만족함을 증명하라.
\end{exercise}

\begin{exercise}
$\charac K=p>0$이고 $L/K$가 유한확장이라 하자. $p\nmid[L:K]$이면 $L/K$가 분리확장임을 증명하라.
\end{exercise}

\begin{exercise}[*]
문제 \thechapter.10의 확장에서
\[
  K_{\mathrm s}=K(\beta)
\]
임을 엄밀히 증명하고
\[
  [K_{\mathrm s}:K]=2,
  \qquad
  [L:K_{\mathrm s}]=p
\]
를 확인하라.
\end{exercise}

\begin{exercise}
$K$가 완전체이고 $L/K$가 대수적 확장이면 $L$도 완전체임을 증명하라.
\end{exercise}

\begin{exercise}[*]
유한확장 $L/K$가 단순확장일 필요충분조건이 중간체를 유한개만 갖는 것임을 다음 순서로 증명하라.
\begin{enumerate}[label=(\alph*)]
  \item $L=K(\alpha)$이면 각 중간체를 $\alpha$의 그 중간체 위 최소다항식과 연결하라.
  \item $K$가 유한체이면 $L/K$가 단순확장임을 보여라.
  \item $K$가 무한하고 중간체가 유한개뿐이면, $L=K(\alpha,\beta)$인 경우 체들 $K(\alpha+c\beta)$를 비교하라.
  \item 생성원 수에 대한 귀납법으로 결론을 완성하라.
\end{enumerate}
\end{exercise}
